% Ad Atticum 14.16 (ad-atticum-14-16) --- English translation
% Translated by Claude (Anthropic) from the Latin (Perseus).
% Style guide: STYLE.md.

\ciceroLetterOpener{CICERO ATTICO salutem}

\ciceroSection{1} I send this letter on the fifth before the Nones, going aboard from the Cluvian gardens into a rowing-pinnace, having handed over Pilia's villa by the Lucrine, with the bailiffs and agents. As for myself, that day I was bearing down upon our Paetus's cheese-and-salt-fish; in a very few days, on to Pompeii; after that, back by water to these realms of Puteoli and the Cumean. O places otherwise much to be wished for --- but with their crowd of interrupters, almost to be fled!

\ciceroSection{2} But to come to the matter --- O the great heroism \textit{[Greek: aristeian]} of our Dolabella! What a reconsideration \textit{[Greek: anathe\=or\=esis]} it is! For my part I do not stop praising him and urging him on. You rightly indicate in every letter what you think of the thing and of the man. Our Brutus, it seems to me, could now even carry a crown of gold through the Forum. For who would dare to harm him, with cross and rock set up before their eyes --- especially amid such applause, such approval from the lowest classes?

\ciceroSection{3} Now, my Attic, see that you set me free. I want, once I have satisfied our Brutus fully enough, to run over to Greece. It matters greatly to my Cicero --- or rather to me, or by Hercules to both of us --- that I should be on hand while he studies. For Leonides' letter that you sent me --- what, I ask, does it contain that we should greatly rejoice in? He will never seem to me sufficiently praised when he is praised in the way he is being praised now. This is not the testimony of a man who trusts him but rather of one who fears for him. I had instructed Herodes to write to me thread by thread \textit{[Greek: kata miton]}. From him so far there is not a single letter. I am afraid he has had nothing that he thought would, on my learning it, give me pleasure.

\ciceroSection{4} What you wrote to Xeno is most welcome to me, for that Cicero want for nothing is a matter both of duty and of my own reputation. I hear that Flamma Flaminius is at Rome. I have written to him that I have instructed you by letter to speak with him about Montanus's business, and I should be glad if you would see to the delivery of the letter I sent him and yourself, when it suits your convenience, talk it over with him. I think, if there is any shame in the man, he will see to it that the payment is not made late and at a loss. About Attica you have done me a great favour by seeing to it that I knew she was well before I knew she had not been quite right.
